\section{路径 B：阅读器多样性与位置剖面重测}

\subsection{B1：阅读器选择（P2）}

单一阅读器是``premature''观感的放大器：所有端到端结论都隐含``对 GLM-4-Plus 成立''的限定。补强目标是把结论升格为``\sieve{} 作为前置选择器的行为具有跨阅读器趋势稳定性''——是趋势级主张，不是逐阅读器显著性主张，样本量与预算都按此设计。阅读器集建议三档：保留 GLM-4-Plus（既有数据的连续性锚点）；增加一个不同家族的旗舰 API 阅读器（例如 GPT-4o-mini 或 DeepSeek-V3 级，取当时可得且计费稳定者）；增加一个本地开源阅读器（Llama-3.1-8B-Instruct 或 Qwen2.5-7B-Instruct，经 vLLM 在单张消费级 GPU 上服务）。本地阅读器有双重价值：不受 API 限流（A0 的阻塞因素不再复现），且完全可复现——审稿人可以一字不差地重跑。

协议与 A 路径严格同构：同一池（QASPER 前 100 条）、同一方法集（七方法）、同一 $r{=}4$、同一提示模板与解码参数（temperature 0、96 token 上限）；唯一变量是阅读器端点。判定标准预先声明：跨三个阅读器，比较 \sieve{} 相对 head 的 F1 保留率差（配对区间）与 bm25--sieve 的差——若三个阅读器上同号且量级同阶，趋势稳定成立；若出现异号，如实报告并讨论阅读器家族与压缩敏感性的交互，这本身即是可发表的结构性发现。预算：每阅读器 700 对调用（API 档两项共 $1.4\times10^3$ 对，本地档仅 GPU 时租）。

\subsection{B2：阅读器侧位置剖面重测（P2，低成本高杠杆）}

局限一节的第三条边界承认：U 形位置先验的依据是借自 \texttt{liu2023lost} 的阅读器侧剖面，而未在自己的阅读器上重测。补法是把论文 M4 位置研究的方向反过来：构造一组合成上下文（从池中抽取、把金证据句人工放置在 5 个相对位置 $0/0.25/0.5/0.75/1$——M4 脚本已实现该放置逻辑，方向相反），分别在三个阅读器上以 full-context 方式提问，记录答案 F1 随证据位置的曲线。每曲线 $5\times n$ 次调用（$n{=}50$ 已足够画出剖面），总预算数百次调用。产出直接回填两处写作：U 形先验的动机从``借用''升格为``自测''；若某阅读器剖面显著非 U 形（例如强中庸偏好），先验的 per-reader 校准即刻成为新实验方向——这是 formulation 框架自然生长出来的问题，恰是第一阶段要求的``问题形式化引领实验''的示范。

\subsection{B3：与相关研究的合并}

阅读器扩展后，召回--F1 相关研究（论文第八节）的样本从``一个阅读器 $\times$ 若干方法''升格为``三个阅读器 $\times$ 若干方法''，\texttt{run\_correlation.py} 的合并逻辑无需改动，只需按阅读器分组重跑。若``弱相关''在三个阅读器上一致成立，``两种协议分开报告、永不合并''的方法论主张获得跨阅读器背书；若相关强度随阅读器显著变化，则获得第三类结构性发现（阅读器对证据召回的敏感度不同），并直接为下一版的讨论节提供素材。
